% ── Abstract ─────────────────────────────────────────────
\chapter*{Abstract}
\addcontentsline{toc}{chapter}{Abstract}
\thispagestyle{plain}

The rapid proliferation of Internet of Things (IoT) devices has placed
growing pressure on mobile edge computing (MEC) infrastructure, where
edge cloudlets must handle large volumes of latency-sensitive tasks
under tight resource constraints. Conventional load balancing methods
that rely on fixed rules or purely reactive heuristics tend to
struggle in such settings because they neither anticipate future
congestion nor adapt their policies based on past experience.

This thesis presents \textbf{DRL-LSTM-LBRO}, a two-stage intelligent
task scheduling framework that addresses these limitations by combining
short-horizon load forecasting with deep reinforcement learning-based
placement decisions. In the first stage, a per-cloudlet GRU-LSTM
predictor estimates near-future CPU and RAM utilisation from a sliding
window of historical observations. These predictions are incorporated
into the state representation seen by a Double Deep Q-Network (DDQN)
agent, which learns a task placement policy through interaction with a
simulated MEC environment modelled as an M/M/$c$ queuing system. The
reward function jointly penalises task drops, deadline violations,
end-to-end latency, energy consumption, and load imbalance, preventing
the well-known policy-collapse behaviour where an agent routes all
traffic to the most capable node.

Simulation-based evaluation on a three-cloudlet topology using task
workloads derived from the Google Borg cluster trace shows that
DRL-LSTM-LBRO achieves a \textbf{0.58\%} average task drop rate and
\textbf{0.260~s} average end-to-end latency over 500 evaluation
episodes (ten random seeds). Compared to five published load balancing
algorithms and a Proximal Policy Optimisation (PPO) baseline, the
proposed framework reduces drop rates by up to 88.9\% and latency by
up to 26.3\%. A scalability experiment extending the topology to six
cloudlets confirms that the framework generalises without architectural
modifications, achieving a 37.4\% drop reduction over the classical
LBRO baseline.

\vspace{0.5cm}
\noindent\textbf{Keywords:} Mobile edge computing, load balancing,
deep reinforcement learning, Double DQN, LSTM, IoT task scheduling,
M/M/$c$ queuing, scalability.
